\chapter{Implementation and Code Organization}
\label{app:implementation_code_organization}

\section{Repository Structure}
\label{sec:appendix_repositories}

The experimental implementation is divided across three repositories. The
\texttt{taming-transformers} repository contains the VQGAN model and its
experiment scripts~\cite{tamingtransformersgithub}, while \texttt{LlamaGen}
contains the corresponding tokenizer implementation~\cite{llamagengithub}.
The \texttt{vq-explore} repository contains the
comparative analysis notebooks and the decoder-locality analysis utility. This
separation keeps model-specific inference code close to each model while
placing cross-model analysis in one location.

\section{Experiment Scripts}
\label{sec:appendix_experiment_scripts}

Equivalent scripts were implemented for the two tokenizers. The reconstruction
scripts encode and decode image-folder datasets and export per-image fidelity
metrics and reconstructed images. The code-usage scripts accumulate global and
position-dependent token counts. The robustness exporters implement the three
interventions used in this thesis: global Gaussian noise, local image-space
patch perturbations, and local token-space replacement.

For each robustness run, images are processed in batches on the GPU. Image
writing and metadata export are delegated to worker processes so that disk
operations do not block inference. Each sample is stored under its dataset
class and image identifier. Sharded JSONL metadata records the perturbation
parameters, patch coordinates, and clean and perturbed token grids. This format
allows later analyses to recompute metrics without rerunning either tokenizer.

The token-space experiments additionally require mappings from each token to a
closest, farthest, and approximately orthogonal alternative. LlamaGen mappings
are computed over its full, L2-normalized codebook. VQGAN mappings are computed
over the empirically active subset used by the experiment, as described in
Chapter~\ref{ch:rq4}.

\section{Analysis Notebooks and Outputs}
\label{sec:appendix_notebook_analysis}

One notebook consolidates the analysis for each experimental chapter:
\texttt{ch04\_rq1\_reconstruction\_codebook\_usage.ipynb},
\texttt{ch05\_rq2\_global\_robustness.ipynb},
\texttt{ch06\_rq3\_local\_image\_perturbations.ipynb},
\texttt{ch07\_rq4\_token\_space\_perturbations.ipynb}, and
\texttt{ch08\_rq5\_distribution\_shift.ipynb}. The notebooks read the
exported images and metadata, compute chapter-level summaries, and save the
tables and figures cited in the thesis. Expensive per-image calculations are
cached as CSV, Parquet, or NPZ files where appropriate. CPU-bound image metrics
are parallelized across workers, while LPIPS is evaluated in GPU batches.

On the analysis host, chapter outputs follow a common layout:
\texttt{chapterN\_outputs/data} contains summary tables,
\texttt{chapterN\_outputs/figures} contains plots, and
\texttt{chapterN\_outputs/qualitative\_samples} contains the source images used
for visual comparisons. Compact source artifacts are copied into
\texttt{Assets/rqN/data} and \texttt{Assets/rqN/figures} in the thesis
repository, where \texttt{rqN} denotes the research-question number. Figures
included by LaTeX are copied into \texttt{Images}. Raw reconstructions and
complete run exports remain outside the thesis repository.

\section{Reproducibility Record}
\label{sec:appendix_reproducibility}

The file \texttt{Assets/reproducibility\_manifest.csv} records the repository
revision, working-tree state, and SHA-256 digest of each experiment script and
analysis notebook in the audited snapshot. The revision identifies the last
commit that changed the file. The digest identifies its exact local contents.
Most entries match their committed revision. The two reconstruction exporters
and the RQ2 notebook contain later uncommitted changes, so the manifest marks
them as modified rather than presenting the listed commit as an exact snapshot.

Both tokenizers process RGB images at $256\times256$ resolution using a
shortest-side resize followed by a center crop. Both expose a nominal
16{,}384-entry codebook and produce a $16\times16$ token grid. The VQGAN runs
load \texttt{/checkpoints/vqgan\_imagenet\_f16\_16384/model.yaml} and
\texttt{last.ckpt}; the LlamaGen runs load
\texttt{/checkpoints/vq\_ds16\_c2i.pt} with the VQ-16 architecture, embedding
dimension 8, and codebook size 16{,}384. Robustness exporters set the Python,
NumPy, and PyTorch random seeds to 0. Code-usage exports are deterministic with
respect to input order: class directories and file names are sorted, data
loading uses \texttt{shuffle=False}, and no token sampling is performed.

For each code-usage run, \texttt{global\_counts.npy} stores one count per
codebook index and \texttt{position\_counts.npy} stores counts with shape
$K\times16\times16$. The exporter also writes \texttt{usage.csv}, the 500 most
frequent codes, and \texttt{summary.json}. When enabled,
\texttt{indices.npz} stores one token grid and source path per image. The
Chapter~\ref{ch:rq5} artifact inventory verifies image counts, token-grid shape,
total token counts, and active-code counts for every model--dataset pair. These
checks are archived in \texttt{Assets/rq5/data/ch08\_artifact\_inventory.csv}.
